\documentclass[10pt,a4paper]{amsart}
\usepackage[margin=19mm]{geometry}
\usepackage{amsmath,amssymb,mathtools}
\usepackage[scr=rsfs,cal=euler]{mathalfa}
\usepackage{xfrac}
\usepackage[unicode,hidelinks,bookmarks=false]{hyperref}
\newcommand{\ie}{i.e.\ }
\newcommand{\cf}{cf.\ }
\newcommand{\resp}{respectively\ }
\newcommand{\Subsection}{\subsection}
\providecommand{\nobreakdash}{\nobreak\hbox{-}}
\setlength{\emergencystretch}{2em}
\multlinegap0pt
\usepackage{bm}
\usepackage{bbm}
\usepackage{dsfont}
\usepackage{xspace}
\usepackage{cancel}
\usepackage{mathtools}
\usepackage{amsthm}
\makeatletter
\@ifundefined{lemma}{\newtheorem{lemma}{Lemma}}{}
\makeatother
\newcommand{\E}{\mathbb E}
\newcommand{\R}{\mathbb{R}}
\newcommand{\cF}{\mathcal{F}}
\newcommand{\norm}[1]{\|#1\|}
\renewcommand{\epsilon}{\varepsilon}
\title[Mathematical methods of reinforcement learning]{Mathematical methods of reinforcement learning}
\author{Denis Belomestny, Alexander Gasnikov, Egor Gladin, Alexey Naumov, Artemy Rubtsov, Yuri Sapronov, Daniil Tiapkin, Nikita Yudin}
\date{Source: arXiv:2607.06935v1 (CC BY 4.0)}
\begin{document}
\maketitle

\noindent\emph{Source and licence.} Mathematical methods of reinforcement learning. Version arXiv:2607.06935v1 (CC BY 4.0), distributed under the Creative Commons Attribution 4.0 International licence, as recorded in the arXiv licence metadata for this exact version and shown on the abstract page https://arxiv.org/abs/2607.06935v1. Full rights notice: licenses/arxiv-2607.06935-CC-BY-4.0.txt.\par\smallskip

\noindent\textit{Editorial scope.} Counted unit: the Lemma whose source label is \texttt{lem:q\_learning\_martingale\_contraction} in the article (source0.tex lines 1224--1264, source label \texttt{lem:q\_learning\_martingale\_contraction}), delivered together with the article's own complete original proof of it (source0.tex lines 1266--1363, one \texttt{proof} environment of 98 lines). No mathematics is written, repaired or re-derived: statement and proof are the article's own text line for line. Required context retained uncounted: none. Every definition, display and label of those lines is retained unchanged. Omitted from this package and not redistributed: 1--1223, 1265--1265, 1364--3047. Nothing inside a retained excerpt is omitted or altered: every retained body is delivered exactly as the article's lines are written, so no omission receipt of any kind accompanies this package. Citations: the retained excerpts carry no citation command at all, so no bibliography entry of the article is transcribed and no reference is redistributed. Reference adaptation: none was needed and none was made; every reference inside the delivered unit resolves inside it, so no unresolved reference or citation is produced. Printed slips kept literal: no printed word, symbol, hypothesis or number of the counted statement or of its proof was altered. Layout and class: the article's own class is \texttt{article} with packages including arxiv\_style, lmodern, tikz, fontenc, inputenc, babel, amsmath, amsfonts, amssymb, mathrsfs, amscd, longtable, amsthm, xy; the wrapper is the lane's pinned \texttt{amsart} editorial wrapper at 10pt on A4, which restates the theorem-like environments the retained text needs and adds the article's own macro definitions that the retained text needs (\texttt{\textbackslash E}, \texttt{\textbackslash R}, \texttt{\textbackslash cF}, \texttt{\textbackslash epsilon}, \texttt{\textbackslash norm}), with extra wrapper packages (bm, bbm, dsfont, xspace, cancel, mathtools). The delivered numbering is the unit's own. Assets: no retained span contains a figure environment, a graphics-inclusion command, a PDF-page inclusion, a table or a TikZ picture; no image file is copied into this package, the package declares no asset dependency and no image file is redistributed with it. Licence: version arXiv:2607.06935v1 of this article is distributed under the Creative Commons Attribution 4.0 International licence, as recorded in the arXiv licence metadata for this exact version and shown on the abstract page https://arxiv.org/abs/2607.06935v1. A full rights notice is delivered at licenses/arxiv-2607.06935-CC-BY-4.0.txt. Characters: every retained excerpt is pure ASCII, so the delivered engine, XeLaTeX with its default Latin Modern text font, reports no missing character.
\begin{lemma}
\label{lem:q_learning_martingale_contraction}
Let \(u_t\in\R^d\) be an adapted bounded process satisfying
\begin{equation}
\label{eq:q_learning_abstract_martingale_recursion}
    u_t
    =
    (1-\alpha_t)u_{t-1}
    +
    \alpha_t(g_{t-1}+\zeta_t),
    \qquad t\in\mathbb{N}.
\end{equation}
Assume that
\[
    \alpha_t\in[0,1],
    \qquad
    \sum_{t=1}^{\infty}\alpha_t=\infty,
    \qquad
    \sum_{t=1}^{\infty}\alpha_t^2<\infty,
\]
that \(g_{t-1}\) is \(\cF_{t-1}\)-measurable and satisfies
\[
    \norm{g_{t-1}}_\infty
    \le
    \beta\norm{u_{t-1}}_\infty+c_{t-1}
    \qquad
    \text{for some }\beta\in[0,1),
\]
where \(c_t\to0\) as \(t\rightarrow\infty\), and that \(\zeta_t\) is a martingale-difference noise with bounded conditional second moment:
\[
    \E[\zeta_t\mid\cF_{t-1}]=0,
    \qquad
    \E[\norm{\zeta_t}_\infty^2\mid\cF_{t-1}]\le C<\infty.
\]
Then
\[
    u_t\to0
    \qquad
    \text{almost surely.}
\]
\end{lemma}

\begin{proof}
Fix a coordinate \(i\in\{1,\ldots,d\}\). The process
\[
    N_n^{(i)}:=\sum_{t=1}^{n}\alpha_t\zeta_t(i)
\]
is a square-integrable martingale whose predictable quadratic variation is bounded, since
\[
    \sum_{t=1}^{\infty}
    \E[\alpha_t^2\zeta_t(i)^2\mid\cF_{t-1}]
    \le
    C\sum_{t=1}^{\infty}\alpha_t^2
    <
    \infty.
\]
Hence the martingale convergence theorem implies that \(N_n^{(i)}\) converges almost surely. Therefore its tails vanish when \(n\rightarrow\infty\):
\begin{equation}
\label{eq:q_learning_unweighted_martingale_tail}
    \sup_{m\ge n}
    \left|
    \sum_{t=n}^{m}\alpha_t\zeta_t(i)
    \right|
    \longrightarrow0
    \qquad
    \text{almost surely.}
\end{equation}
For integers \(a\le b\), define
\[
    \Phi_{a:b}:=\prod_{k=a}^{b}(1-\alpha_k),
    \qquad
    \Phi_{a:b}:=1\quad\text{if }a>b.
\]
Unrolling~\eqref{eq:q_learning_abstract_martingale_recursion} from time \(n\) to time \(m\) yields
\begin{equation}
\label{eq:q_learning_unrolled_abstract_recursion}
    u_m(i)
    =
    \Phi_{n+1:m}u_n(i)
    +
    \sum_{t=n+1}^{m}\Phi_{t+1:m}\alpha_t g_{t-1}(i)
    +
    \sum_{t=n+1}^{m}\Phi_{t+1:m}\alpha_t\zeta_t(i).
\end{equation}
By summation by parts, because the weights \(\Phi_{t+1:m}\) lie in \([0,1]\) and are monotone in \(t\), the weighted martingale tails are controlled by the unweighted tails. In particular, when \(n\rightarrow\infty\)
\begin{equation}
\label{eq:q_learning_weighted_martingale_tail}
    \sup_{m\ge n}
    \left|
    \sum_{t=n+1}^{m}\Phi_{t+1:m}\alpha_t\zeta_t(i)
    \right|
    \longrightarrow0
    \qquad
    \text{almost surely.}
\end{equation}
Let
\[
    L:=\limsup_{t\to\infty}\norm{u_t}_\infty.
\]
Since \(u_t\) is bounded, \(L<\infty\). Fix \(\epsilon>0\). Choose \(n\) large enough so that, for all \(t\ge n\),
\[
    \norm{u_t}_\infty\le L+\epsilon,
    \qquad
    c_t\le\epsilon,
\]
and so that the weighted martingale tail in~\eqref{eq:q_learning_weighted_martingale_tail} is at most \(\epsilon\) for every coordinate. Using~\eqref{eq:q_learning_unrolled_abstract_recursion}, we obtain, for every \(m\ge n\),
\begin{align*}
    |u_m(i)|
    &\le
    \Phi_{n+1:m}|u_n(i)|
    +
    \sum_{t=n+1}^{m}\Phi_{t+1:m}\alpha_t
    \left(\beta\norm{u_{t-1}}_\infty+c_{t-1}\right)
    +
    \epsilon
    \\
    &\le
    \Phi_{n+1:m}|u_n(i)|
    +
    \left(1-\Phi_{n+1:m}\right)
    \left(\beta(L+\epsilon)+\epsilon\right)
    +
    \epsilon,
\end{align*}
where we used
\[
    \sum_{t=n+1}^{m}\Phi_{t+1:m}\alpha_t
    =
    1-\Phi_{n+1:m}.
\]
Since \(\sum\limits_{t = 1}^{\infty}\alpha_t=\infty\), we have \(\Phi_{n+1:m}\to0\) as \(m\to\infty\). Taking \(\limsup\limits_{m\to\infty}\) and maximizing over the finitely many coordinates gives
\[
    L\le \beta(L+\epsilon)+2\epsilon.
\]
Equivalently,
\[
    (1-\beta)L\le(\beta+2)\epsilon.
\]
Letting \(\epsilon\rightarrow0\) yields \(L=0\), because \(\beta<1\). Thus \(u_t\to0\) when \(t\rightarrow\infty\) almost surely. This concludes the lemma proof.
\end{proof}

\end{document}
